%%%%%%%%%%%%%%%%%%%%%%%%%%%%%%%%%%%%%%%%%%%%%%%%%%%%%%%%%%%%%%%%%%%%%%%%%%%%%%%%%%%%%%%%%%%%%%%%%%%%%%%%%%%%%%%%%%%%%%%%
\section{Project Initiation}
\label{sec:init}
%%%%%%%%%%%%%%%%%%%%%%%%%%%%%%%%%%%%%%%%%%%%%%%%%%%%%%%%%%%%%%%%%%%%%%%%%%%%%%%%%%%%%%%%%%%%%%%%%%%%%%%%%%%%%%%%%%%%%%%%
The project initiation phase begins once a project has been approved for funding. Following the communication of the 
awarded project by the Program Officer (PO), the responsible Section Head (SH) initiates the project setup. The purpose
of this phase is to establish the project within the eScience Center by completing the required administrative setup,
assigning responsibilities, and preparing the project for execution.

The exact initiation process depends on the project type. External projects, as well as certain internal projects, may
require additional contractual or legal arrangements before they can commence. Further information on the different 
project types is provided in Appendix~\ref{app:internal-projects} and Appendix~\ref{app:external-projects}.

Once the project has been registered, staffed, and all applicable administrative, financial, and legal requirements have 
been fulfilled, it is ready to enter the execution phase.

%%%%%%%%%%%%%%%%%%%%%%%%%%%%%%%%%%%%%%%%%%%%%%%%%%%%%%%%%%%%%%%%%%%%%%%%%%%%%%%%%%%%%%%%%%%%%%%%%%%%%%%%%%%%%%%%%%%%%%%%
\subsection{SH assignment}
\label{sec:init:pm-assign}
%%%%%%%%%%%%%%%%%%%%%%%%%%%%%%%%%%%%%%%%%%%%%%%%%%%%%%%%%%%%%%%%%%%%%%%%%%%%%%%%%%%%%%%%%%%%%%%%%%%%%%%%%%%%%%%%%%%%%%%%
Each project has one assigned SH, who is accountable for the project throughout its lifecycle.

For projects funded through the eScience Center’s regular funding calls, the assigned Section Head normally corresponds 
to the section under which the proposal was submitted. In other cases, or where the assignment is not self-evident, the
Section Heads jointly determine the most appropriate assignment, taking into account the project’s scope, expertise,
and organizational responsibilities.

Once the assignment has been agreed, the responsible Section Head is recorded and Finance is notified to initiate the
project registration process.

\textbf{Responsible: Section Heads.}

%%%%%%%%%%%%%%%%%%%%%%%%%%%%%%%%%%%%%%%%%%%%%%%%%%%%%%%%%%%%%%%%%%%%%%%%%%%%%%%%%%%%%%%%%%%%%%%%%%%%%%%%%%%%%%%%%%%%%%%%
\subsection{Project registration}
\label{sec:init:registration}
%%%%%%%%%%%%%%%%%%%%%%%%%%%%%%%%%%%%%%%%%%%%%%%%%%%%%%%%%%%%%%%%%%%%%%%%%%%%%%%%%%%%%%%%%%%%%%%%%%%%%%%%%%%%%%%%%%%%%%%%
The table below summarizes the main project registration outputs, the responsible stakeholder, and the associated
activities.

\renewcommand{\TblrNewPage}{\clearpage}  %% <--- Added this line only
\begin{longtblr}[theme = fancy,
    caption= {Project registration activities and responsibilities},
     label = {tab:init:registration},
]{
  colspec={|p{0.15\textwidth}|p{0.15\textwidth}|p{0.65\textwidth}|}, width = 1\linewidth,
  rowhead = 1, %rowfoot = 1,
  row{1} = {font=\bfseries},
  column{1} = {font=\bfseries},
  row{odd} = {white}, row{even} = {white},
  row{1} = {gray!20},
  cell{4}{1} = {r=2}{t},
}
\toprule
Output & Responsible & Main actions \\
\toprule
Project record & Finance &
\begin{minipage}[t]{1\linewidth}
\begin{itemize}\itemsep0em
    \item Creating the project code in AFAS
    \item Assigning the responsible Section Head as budget holder
\end{itemize}
\end{minipage} \\
\midrule
Project portfolio & Finance &
\begin{minipage}[t]{1\linewidth}
\begin{itemize}\itemsep0em
    \item Creating the project portfolio folder~\cite{proj-portfolio} following the rquired template subfolders
    \item Requesting missing project documentation
    \item     Uploading the grant package (proposal, awarding letter, signed start form, and supporting documentation
    submitted by the LA)
\end{itemize}
\end{minipage} \\
\midrule
\end{longtblr}
\textbf{Responsible: Finance.}

%%%%%%%%%%%%%%%%%%%%%%%%%%%%%%%%%%%%%%%%%%%%%%%%%%%%%%%%%%%%%%%%%%%%%%%%%%%%%%%%%%%%%%%%%%%%%%%%%%%%%%%%%%%%%%%%%%%%%%%%
\subsection{Project staffing}
\label{sec:init:staffing}
%%%%%%%%%%%%%%%%%%%%%%%%%%%%%%%%%%%%%%%%%%%%%%%%%%%%%%%%%%%%%%%%%%%%%%%%%%%%%%%%%%%%%%%%%%%%%%%%%%%%%%%%%%%%%%%%%%%%%%%%
The responsible SH is responsible for staffing the project, including the appointment of a Lead RSE and the assembly of 
the eScience Center project team~\cite{rse-teams}. Staffing decisions take into account the project’s requirements, 
available capacity, and the expertise required for successful delivery.
Staffing may be adjusted throughout the project as needed (see Section~\ref{sec:exec:changes}). The Lead RSE serves as 
the primary point of contact for the LA, unless otherwise agreed.

%%%%%%%%%%%%%%
\subsubsection{Staffing preparation}
%%%%%%%%%%%%%%
The responsible Section Head reviews the project proposal to determine the staffing requirements for the project. The 
assessment considers the project scope and objectives, the awarded budget, the expected timeline, the required 
scientific and technical expertise (making use of available organizational information on technology capabilities, e.g., 
technology skills surveys), any applicable contractual, legal, or organizational constraints, and the available capacity 
within the section based on the organizational planning system. These considerations form the basis for assigning the 
eScience Center project team.

%%%%%%%%%%%%%%
\subsubsection{Team assignament}
%%%%%%%%%%%%%%
Every project must have: one accountable Lead RSE; and at least one additional Research Software Engineer assigned to 
the project.
The eScience Center project team is assigned according to the following process:
\begin{enumerate}
    \item The responsible SH identifies a suitable candidate for the Lead RSE role, taking into account the staffing 
    preparation described above.
    \item The SH discusses the proposed assignment with the candidate.
    \item The proposed staffing is reviewed and confirmed by the Section Heads.
    \item The responsible SH formally assigns the Lead RSE and the remaining project team members.
\end{enumerate}

%%%%%%%%%%%%%%
\subsubsection{Project planning}
%%%%%%%%%%%%%%
Following project registration, the project is synchronized to the organization’s planning system. The responsible SH 
records the Lead RSE in the corresponding project in Ganttic. The Lead RSE prepares the initial project planning by 
allocating the assigned team members across the different project phases. The planning serves as the baseline for 
project execution and may be updated during the project as required.

\textbf{Responsible: Section Head.}

%%%%%%%%%%%%%%%%%%%%%%%%%%%%%%%%%%%%%%%%%%%%%%%%%%%%%%%%%%%%%%%%%%%%%%%%%%%%%%%%%%%%%%%%%%%%%%%%%%%%%%%%%%%%%%%%%%%%%%%%
\subsection{Administrative kick-off meeting}
\label{sec:init:admin-kickoff}
%%%%%%%%%%%%%%%%%%%%%%%%%%%%%%%%%%%%%%%%%%%%%%%%%%%%%%%%%%%%%%%%%%%%%%%%%%%%%%%%%%%%%%%%%%%%%%%%%%%%%%%%%%%%%%%%%%%%%%%%
Once the project has been registered and staffed, the responsible Section Head (SH) organizes an administrative kick-off 
meeting with the Lead Applicant (LA). Prior to the meeting, the SH and Lead RSE review the project proposal and 
identify any project-specific requirements that should be discussed with the LA. The purpose of the meeting is to 
introduce the collaboration with the Netherlands eScience Center and to ensure that the project is administratively and 
organizationally ready to start. These may include:
\begin{itemize}
\item any additional organizational or technical support required for the project;
\item deviations from the default intellectual property or software licensing arrangements;
\item data readiness, data access requirements, and any personal data or GDPR requirements;
\item software sustainability expectations; and
\item ethical, legal, or other project-specific risks.
\end{itemize}

The meeting is supported by a standardized administrative kick-off presentation available in the project portfolio under 
the \textit{Templates projects} folder~\cite{proj-portfolio}. The presentation is periodically updated to reflect the
requirements of different funding calls and the eScience Center's current procedures and policies. The responsible
SH uses this presentation to ensure that all standard topics are covered consistently across projects while adapting
the discussion to the specific needs of the project.

The presentation typically covers the following topics:
\begin{itemize}
\item Introduction to the Netherlands eScience Center (mission, organization, and expertise);
\item Working with the eScience Center (roles and responsibilities, collaboration model, in-kind and in-cash 
contributions); 
\item Project lifecycle and project governance;
\item Community activities and research impact (Research Software Directory, Digital Skills Programme, 
Fellowship Programme, etc.);
\item Intellectual property, software licensing, publication practices, and open science principles;
\item Project-specific needs, including scientific and technical expertise, SURF infrastructure requirements,
 and additional resources (e.g., AI subscriptions or other project-specific services).
\end{itemize}

The presentation used during the meeting and any meeting notes are stored in the corresponding project portfolio folder
together with the other project documentation. The Secretary may assist the Section Head with scheduling the meeting.

\begin{table}[!h]
\captionsetup{
  labelformat=fancytable,
  labelsep=space,
  labelfont=bf,
  textfont=bf,
  justification=centering,
  singlelinecheck=true
}
\caption{Administrative kick-off meeting overview}
\centering 
\begin{booktabs}{
   colspec={|>{\bfseries}p{0.16\textwidth}|p{0.74\textwidth}|},
   row{even}={gray!20}
 }
 \toprule
 Timing: & Soon after project registration and staffing have been completed. \\[1.2ex]
 Participants: &  SH (organizer), LA, Lead RSE (optional). Others may be invited as needed. \\[1.2ex]
 Objective: &  Introduce the collaboration with the Netherlands eScience Center, discuss the project setup, and identify 
 any project-specific administrative or organizational requirements before project execution begins. \\[1.2ex]
 Duration: &  1 hour. \\[1.2ex]
 Location: &  Online. \\[1.2ex]
 \bottomrule
 \end{booktabs}
\end{table}

\textbf{Responsible: Section Head.}

%%%%%%%%%%%%%%%%%%%%%%%%%%%%%%%%%%%%%%%%%%%%%%%%%%%%%%%%%%%%%%%%%%%%%%%%%%%%%%%%%%%%%%%%%%%%%%%%%%%%%%%%%%%%%%%%%%%%%%%%
\subsection{Kick-off meeting}
\label{sec:init:kickoff}
%%%%%%%%%%%%%%%%%%%%%%%%%%%%%%%%%%%%%%%%%%%%%%%%%%%%%%%%%%%%%%%%%%%%%%%%%%%%%%%%%%%%%%%%%%%%%%%%%%%%%%%%%%%%%%%%%%%%%%%%
Following the administrative kick-off meeting, the Lead RSE organizes the project kick-off meeting with the 
project team, comprising the eScience project team, the Lead Applicant (LA), and the LA's team. The purpose of the 
meeting is to establish a shared understanding of the project objectives, work plan, responsibilities, and expected 
deliverables before project execution begins.

Prior to the meeting, and based on their review of the project proposal, the Lead RSE and the eScience Center project 
team identify any scientific or technical aspects that should be discussed with the LA. These may include:

\begin{itemize}
\item the scientific objectives and expected project outcomes;
\item the proposed work packages, milestones, timeline, and overall feasibility;
\item technical challenges, assumptions, implementation risks, and whether additional technical expertise or 
consultation from other RSEs is required; 
\item software sustainability considerations and the initial Software Management Plan (SMP);
\item dependencies on data, infrastructure, or external collaborators; and
\item any issues requiring clarification before project execution starts.
\end{itemize}

The meeting is supported by a standardized project kick-off agenda available in the project portfolio under the 
\textit{Templates projects} folder~\cite{proj-portfolio}. The Lead RSE uses this agenda to ensure that all standard 
topics are covered consistently across projects, while adapting the discussion to the specific needs of the project.

The agenda typically covers the following topics:

\begin{itemize}
\item Round of introductions;
\item Project introduction and scientific objectives;
\item Review of the project work plan;
\item Agreements on the initial project planning and deliverables;
\item Software sustainability and the Software Management Plan (SMP);
\item Any other business.
\end{itemize}

The meeting outcomes, agreed action items, and any presentation material are stored in the corresponding project 
portfolio folder. The Lead RSE updates the project planning accordingly and prepares the project for the execution 
phase. These documents establish the initial project baseline and provide the starting point for the project log 
maintained during the execution phase. The Secretary may assist the Lead RSE with scheduling the meeting.

\begin{table}[!h]
\captionsetup{
  labelformat=fancytable,
  labelsep=space,
  labelfont=bf,
  textfont=bf,
  justification=centering,
  singlelinecheck=true
}
\caption{Project kick-off meeting overview}
\centering
\begin{booktabs}{
  colspec={|>{\bfseries}p{0.16\textwidth}|p{0.74\textwidth}|},
  row{even}={gray!20}
}
\toprule
Timing: &
Shortly after the administrative kick-off meeting. \\[1.2ex]
Participants: &
Lead RSE (organizer), SH, LA, project team. Others may be invited as needed. \\[1.2ex]
Objective: &
Establish a shared understanding of the scientific objectives, work plan, responsibilities, and expected deliverables 
before project execution begins. \\[1.2ex]
Duration: &
1.5 hours. \\[1.2ex]
Location: &
At the eScience Center or the project location.\\[1.2ex]
\bottomrule
\end{booktabs}
\end{table}
\smallskip
\textbf{Note:} The Lead RSE, responsible SH, and LA are expected to attend the meeting in person at the eScience Center 
or the project location. Other project team members may participate online if appropriate.

\textbf{Responsible: Lead RSE.} 

%%%%%%%%%%%%%%%%%%%%%%%%%%%%%%%%%%%%%%%%%%%%%%%%%%%%%%%%%%%%%%%%%%%%%%%%%%%%%%%%%%%%%%%%%%%%%%%%%%%%%%%%%%%%%%%%%%%%%%%%
\subsection{Legal agreements}
\label{sec:init:legal}
%%%%%%%%%%%%%%%%%%%%%%%%%%%%%%%%%%%%%%%%%%%%%%%%%%%%%%%%%%%%%%%%%%%%%%%%%%%%%%%%%%%%%%%%%%%%%%%%%%%%%%%%%%%%%%%%%%%%%%%%

The eScience Center promotes open, reproducible science and open-source software development. Projects funded through 
the eScience Center are carried out under the eScience Center Terms and Conditions~\cite{nlesc-terms}. However, project 
partners may require additional legal agreements to grant access to their facilities, data, infrastructure, or other 
resources. These agreements may also contain provisions related to intellectual property, confidentiality, data 
protection, or information security.

RSEs must not sign any formal agreements without consulting the responsible SH. Such documents include, but are not 
limited to:

\begin{itemize}
\item Guest agreement
\item Data sharing agreement
\item Non-disclosure agreement
\item Consortium agreement
\item Collaboration agreement
\item Intellectual property (IP) agreement
\end{itemize}

The SH reviews the proposed agreement, consulting the Lead RSE where project-specific input is required, and determines 
whether additional review is required. If needed, the SH consults the Policy \& Compliance Officer (P\&CO) before the 
agreement is signed. The final signed agreement is archived in the corresponding project portfolio.

\textbf{Responsible: Section Head.} 

%%%%%%%%%%%%%%%%%%%%%%%%%%%%%%%%%%%%%%%%%%%%%%%%%%%%%%%%%%%%%%%%%%%%%%%%%%%%%%%%%%%%%%%%%%%%%%%%%%%%%%%%%%%%%%%%%%%%%%%%
\subsection{Technology plan}
\label{sec:init:techplan}
%%%%%%%%%%%%%%%%%%%%%%%%%%%%%%%%%%%%%%%%%%%%%%%%%%%%%%%%%%%%%%%%%%%%%%%%%%%%%%%%%%%%%%%%%%%%%%%%%%%%%%%%%%%%%%%%%%%%%%%%

Every research project starts with a review of existing literature and technologies, and an eScience project is no exception.
To that end, the Lead RSE should write down a technology plan for the project before executing it. The starting point of 
a technology plan is the eScience section, the workplan and the timeline of the project proposal as well as agreements from the kick-off meeting.

The project team submits a technology plan by the date agreed during the project kick-off, describing
\begin{itemize}
\item available technology options, selected tools for the project, and the rationale for their use
\item the technological outcomes of the project (software and data) 
\item steps to be taken regarding reusability and adoptability, etc. 
\end{itemize}
The technology plan outlines key choices such as programming languages and quality standards. It serves as a living document, 
capturing technical decisions and ensuring optimal use of in-house expertise; see an example in ~\cite{spaaks2023asreview}.

\let\myhcolw\relax 
\newlength{\myhcolw}
\setlength{\myhcolw}{0.8\textwidth}
\begin{table}[!h]
\begin{booktabs}{colspec={|>{\bfseries}m{0.15\textwidth}|m{\myhcolw}|},row{even}={gray!20}}
    \toprule
    Written by: &  Lead RSE, in collaboration with project team (including LA, TL), CMs, and others RSEs or colleagues (e.g. with relevant expertise on the subject), or relevant SIG. \\[1.5ex]
    Target audience: & Project team, TLs, PMs  \\[1.5ex]
    Schedule: &  %
    \begin{minipage}[t]{\myhcolw}
    \begin{itemize}\itemsep0em
        \item written at the start of the project work, before any software development starts,
        \item submitted to PM/TL by email before the deadline agreed during the project kick-off, 
        \item as a part of the project log (either full document in the log or a URL to it). 
    \end{itemize} 
      \end{minipage}
    \\[1.5ex]
    Approved by: & PM after due consultation of TL. \\[1.5ex]
    \bottomrule
\end{booktabs}
\end{table}

The Lead RSE is encouraged to consult colleagues with relevant expertise when drafting the technology plan. CMs can advise 
on engaging the target audience for software reuse and adoption. Since TLs are accountable for safeguarding the suitable technology in the project, the involvement of the
TL in writing the technology plan is important. Therefore, the PM must consult the TL on the technology plan, submitted
by the Lead RSE before any technological decisions are made in the project.

Once the TL approves the technology plan (via email), the project team updates the management plans if needed. The Lead 
RSE logs key decisions (see Section~\ref{sec:exec:log}) and, if necessary, archives related emails in the project portfolio. 
Any updates to the plan during the project should be appended to the original, with an explanation of the changes. 
This ensures continuity, documents lessons learned, and supports knowledge transfer if team changes occur. Plan updates are reviewed during the status update meetings (Section~\ref{sec:exec:status}).
